

\illustration{0.3}{khafre}{Khafre's pyramid at Giza. At the top of the pyramid you can see many remaining casing stones. Image by Daaaveee, \href{https://commons.wikimedia.org/wiki/File:Chichen\_Itza\_3.jpg}{commons.wikimedia.org}}

Friendly but competing archaeologists Ameena and Bilal are searching for the lost entrance to the \emph{Nomarch's Celestial Pyramid of Cauris}.
Over the Millennia, much of the outer layer has been stripped away, so the entrance must lie behind one of the remaining ``casing'' stones -- the smooth limestone blocks that once covered the entire monument -- on the north face of the pyramid.

To avoid interference between their sensor arrays, Ameena and Bilal will divide the north face into two connected regions and scan one region each. 
For fairness, each region must contain the same number of casing stones.

They model the pyramid's north face as a triangular grid.
The grid has $N$ rows and $2N - 1$ columns, with the pyramid occupying a subset of these cells: the middle cell of the first row, the middle three cells of the second row, the middle five of the third row, and so on. 
A nonempty subset of the cells in the pyramid are casing stones.

Your task is to partition the pyramid into two regions $A$ and $B$ such that:
\begin{itemize}
\item Every cell in the pyramid belongs to exactly one region.
\item Each region is connected.
  (A region is connected if it is possible to walk between any two of its cells by moving between adjacent cells.
    Two cells are adjacent if they share a side.)
\item The two regions contain the same number of casing stones.
\end{itemize}


